\question
После изучения древних руин в пустыне археологи обнаружили систему туннелей, соединяющих пять храмов (A, B, C, D, E). Известно, что туннели образуют неориентированный взвешенный граф и имеют следующие длины: A-B 3 км, A-C 5 км, B-D 4 км, C-D 2 км, C-E 6 км, D-E 1 км. Необходимо найти кратчайший путь из храма A в храм E, указать его общую длину, а затем закодировать найденный маршрут в двоичном виде. 

Для кодирования используйте фиксированные коды храмов: A = 000, B = 001, C = 010, D = 011, E = 100. Длину каждого туннеля записывайте как 3-битное двоичное число. Маршрут записывается в следующем порядке: сначала код стартового храма, затем 3  бита длины туннеля и 3 бита кода следующего храма. В ответе укажите кратчайший путь, его суммарную длину и итоговую двоичную последовательность маршрута. Если существует несколько кратчайших путей, нужно выбрать путь, у которого последовательность вершин лексикографически меньше.

---------------

Автор -- Никита Глейм
